\section{Introduction}

Deep learning has achieved impressive successes in various real-world applications, such as computer vision \citep{krizhevsky2012imagenet, goodfellow2014generative, Long_2015_CVPR}, natural language processing \citep{graves2013speech, bahdanau2014neural, young2018recent}, and robotics \citep{gu2017deep}. One notable example of this is in the field of image classification, where the winner of the 2017 ImageNet challenge achieved a top-5 error rate of just 2.25\% \citep{hu2018squeeze} using a training dataset of 1 million labeled high resolution images in 1000 categories. Deep neural networks have been shown to outperform humans in speech recognition, with a 5.15\% word error rate using the LibriSpeech training corpus \citep{panayotov2015librispeech}, which consists of approximately 1000 hours of 16kHz read English speech from 8000 audio books.

The remarkable successes of deep learning have challenged conventional machine learning theory, particularly when it comes to high-dimensional data. Existing literature has established a minimax lower bound of sample complexity $n\gtrsim\epsilon^{-(2s+d)/s}$ for learning $s$-H\"{o}lder functions in $\RR^d$ with accuracy $\epsilon$ \citep{gyorfi2006distribution}. 
This minimax lower bound, however, is far beyond the practical limits. For instance, the images in the ImageNet challenge are of the resolution $224\times224=50176$, while the sample size of $1.2$ million is significantly smaller than the theoretical bound.

%Fortunately, real-world data sets often exhibit low-dimensional structures. Many images consist of projections of a three-dimensional object followed by some transformations, such as rotation, translation, and skeleton. This generating mechanism induces a small number of intrinsic parameters\cite{hinton2006reducing, osher2017low}. More broadly, visual, acoustic, textual, and many other types of data often have low-dimensional geometric structures due to rich local regularities, global symmetries, repetitive patterns, or redundant sampling \cite{tenenbaum2000global, roweis2000nonlinear,coifman2005geometric}. 



Several recent results have attempted to explain the successes of deep neural networks
%such observations by considering nonparametric regression 
by taking the low-dimensional structures of data into consideration\citep{chen2019efficient,chen2022nonparametric,nakada2020adaptive,liu2021besov,schmidt2019deep}. Specifically, \citet{chen2022nonparametric} shows that when the input data are supported on a $p$-dimensional Riemannian manifold embedded in $\RR^d$, deep neural networks can capture the  low-dimensional  intrinsic structures of the manifold. The sample complexity in \citet{chen2022nonparametric} depends on the intrinsic dimension $p$, which circumvents the curse of ambient dimension $d$; \citet{nakada2020adaptive} assumes that the input data are supported on a subset of $\RR^d$ with Minkowski dimension $p$, and establishes a sample complexity similar to \citet{chen2022nonparametric}. \citet{liu2021besov} considers a  classification problem, and show that convolutional residual networks enjoy similar theoretical properties  to \citet{chen2022nonparametric}.

Considering the complexity of real world applications, however, the assumptions of data lying exactly on a low-dimensional manifold or a set with low Minkowski dimension are stringent. To bridge such a gap between theory and practice, we consider a relaxed assumption that the input data $X$ are approximately supported on a subset of $\RR^d$ with certain low-dimensional structures denoted by $\cS$. Roughly speaking, there exists a sufficiently small $\tau$ such that we have $\PP(X\notin\cS)=\tau$, where $\cS$ can be characterized by a new complexity notation -- effective Minkowski dimension. We then prove that under proper conditions, the sample complexity of nonparametric regression using deep neural networks only depends on the effective Minkowski dimension of $\cS$ denoted by $p$. Our assumption arises from practical motivations: The distributions of real-world data sets often exhibit a varying density. In practice, the low-density region can be neglected, if our goal is to minimize the $L_2$ prediction error in expectation. 

Furthermore, we illustrate our theoretical findings by considering nonparametric regression with an anisotropic multivariate Gaussian randomly sampled from $N(0,\Sigma)$ design in $\RR^d$. Specifically, we prove that when the eigenvalues of $\Sigma$ have an exponential decay, we can properly construct $\cS$ with the effective Minkowski dimension $p=\min(\cO(\sqrt{\log n}),d)$. Moreover, when the eigenvalues of $\Sigma$ have a polynomial decay, we can properly construct $\cS$ with the effective Minkowski dimension $p=\min(\cO(n^\gamma,d))$, where $\gamma\in(0,1)$ is a small constant. Our proposed effective Minkovski dimension is a non-trivial generalization of the manifold intrinsic dimension \citep{chen2022nonparametric} or the Minkowski dimension  \citep{nakada2020adaptive}, as both the intrinsic dimension or Minkowski dimension of the aforementioned $\cS$'s are $d$, which can be significantly larger than $p$ for moderate sample size $n$.

An ingredient in our analysis is an approximation theory of deep ReLU networks for $\beta$-Hölder functions \citep{YAROTSKY2017103,nakada2020adaptive,chen2019efficient}. Specifically, we show that, in order to uniformly approximate $\beta$-Hölder functions on a properly selected $\cS$ up to an $\epsilon$ error, %which is comparable to the statistical error of nonparametric regression, 
the network consists of at most $O(\epsilon^{-p/\beta})$ neurons and weight parameters, where $p$ is the effective Minkowski dimension of the input data distribution. The network size in our theory only weakly depends on the ambient dimension $d$, which circumvents the curse of dimensionality for function approximation using deep ReLU networks. Our approximation theory is established for the $L^2$ norm instead of the  $L^\infty$ norm in \citet{nakada2020adaptive,chen2019efficient}. The benefit is that we only need to approximate the function accurately on the high-density region, and allow for rough approximations on the low-density region. Such flexibility is characterized by our effective Minkowski dimension.


%which also involve function approximation theory of deep ReLU networks: Chen et al. 2021 and Nakada et al. 2021 requires deep ReLU networks to achieve an accurate point-wise approximation to the target function on the entire support, while our theory only requires deep ReLU networks to approximate the function within a properly selected subset of the support.

The rest of this paper is organized as follows: Section 2 reviews the background; Section 3 presents our functional approximation and statistical theories; Section 4 provides an application to Gaussian random design; Section 5 presents the proof sketch of our main results; Section 6 discusses related works and draws a brief conclusion.

\paragraph{Notations} Given a vector $v=(v_1,...,v_d)^\top\in\RR^d$, we define $\norm{v}_p^p=\sum_{j}|v_j|^p$ for $p\in[1,\infty)$ and $\norm{v}_{\infty}=\max_j|v_j|$. Given a matrix $W=[W_{ij}]\in\RR^{n\times m}$, we define $\norm{W}_{\infty}=\max_{i,j}|W_{ij}|$. We define the number of nonzero entries of $v$ and $W$ as $\norm{v}_0$ and $\norm{W}_0$, respectively. For a function $f(x)$, where $x\in\cX\subseteq \RR^d$, we define $\norm{f}_{\infty} = \max_{x\in\cX}|f(x)|$. We define $\norm{f}^2_{L^2(P)}=\int_{\cX}f^2(x)p(x)dx$, where $P$ is a continuous distribution defined on $\cX$ with the pdf $p(x)$.


%\fbox{Put a main theorem here?}